% from ega4-21c-fr.tex
\oldpage[IV-4]{332}
définition (21.15.2) et de l'exactitude à droite du foncteur $p^*$, que
l'homomorphisme précédent définit par passage aux quotients un homomorphisme canonique
\[
\label{IV.21.15.9.1}
  \operatorname{Div}(X/S)\longrightarrow\operatorname{Div}(X'/S')
  \tag{21.15.9.1}
\]
qui transforme évidemment les éléments de $\operatorname{Div}^+(X/S)$ en éléments de
$\operatorname{Div}^+(X'/S')$. On pose encore
$\operatorname{Div}(X'/S')=\operatorname{Div}_{X/S}(S')$ (resp.
$\operatorname{Div}^+(X'/S')=\operatorname{Div}_{X/S}^+(S')$), et on voit aussitôt
qu'on a ainsi défini deux foncteurs contravariants
\[
  \operatorname{Div}_{X/S}:\mathbf{Sch}_{/S}\longrightarrow\mathbf{Ab},\qquad
  \operatorname{Div}_{X/S}^+:\mathbf{Sch}_{/S}\longrightarrow\mathbf{Ens}
\]
de la catégorie des $S$-préschémas dans celle des groupes commutatifs (resp. des
ensembles). On verra plus loin (chap.~VI) des cas importants où le foncteur
$\operatorname{Div}_{X/S}^+$ est \emph{représentable} $(\mathbf 0_{\mathrm{III}},
8.1.8)$.

Pour tout diviseur $D\in\operatorname{Div}(X/S)$, l'image de $D$ par
l'homomorphisme (21.15.9.1) n'est autre d'ailleurs que l'\emph{image réciproque}
$p^*(D)$ (au sens de (21.4.2)) : l'existence de cette image réciproque et l'assertion
précédente sont en effet des conséquences immédiates de (20.6.5, (iii)) et (20.6.9).

Les éléments de $\operatorname{Div}(X'/S')$ s'appellent souvent, par abus de langage,
« \emph{familles de diviseurs sur $X$ relatifs à $S$, paramétrées par le
$S$-préschéma $S'$} »; on utilise surtout cette terminologie lorsqu'il s'agit de
diviseurs \emph{positifs}.
\end{env}
